\documentclass[noamsfonts, reqno]{amsart}
\usepackage[T1]{fontenc}
\usepackage[utf8]{inputenc}
\usepackage{mlmodern, microtype, mleftright}
\usepackage{mathtools, thmtools, bookmark, cite, amssymb}
\usepackage[nameinlink]{cleveref}
\usepackage[dvipsnames]{xcolor}
\usepackage[ocgcolorlinks]{ocgx2}
\usepackage[bb=boondox, scr=boondox, bfbb]{mathalfa}
\usepackage[margin={3cm, 2cm}, a4paper]{geometry}
\usepackage[foot]{amsaddr}
\hypersetup{linkcolor=BrickRed, citecolor=RoyalBlue}

\usepackage{tikz-cd}

\usepackage[mathlines]{lineno}

\linenumbers
\mleftright

%% ENVIRONMENTS
\declaretheorem[numberwithin=section]{proposition}

%% SHORTCUTS
\newcommand{\catpl}{\mathsf{PreLie}}
\newcommand{\prelie}{\mathbin{\vartriangleright}}
\newcommand{\preliealt}{\mathbin{\blacktriangleright}}
\newcommand{\laction}{\mathbin{\vartriangleright_{\mathrm{L}}}}
\newcommand{\raction}{\mathbin{\vartriangleleft_{\mathrm{R}}}}

\begin{document}
\title[Types of crossed modules]{A survey on types of crossed modules}
\author[J.~Diehl]{Joscha Diehl}
\address{University of Greifswald}
\author[K.~Ebrahimi-Fard]{Kurusch Ebrahimi-Fard}
\address{Norwegian University of Science and Technology}
\author[N.~Tapia]{Nikolas Tapia}
\address{Weierstrass Institute}
\address{HU Berlin}
\date{March 30, 2025}

\maketitle

\section{Introduction}
\label{s:intro}

\section{Crossed modules of pre-Lie algebras}
We take the point of view of \cite{LRW2024}.
In this framework, a pre-Lie crossed module is a pre-Lie structure on a two term chain complex \((\partial\colon T_1\to T_0)\).
This means we have a chain map \(\prelie\colon(T_1\to T_0)\otimes(T_1\to T_0)\to(T_1\to T_0)\) whose associator
\[
  \mathrm{a}_{\prelie}\coloneqq\prelie\circ(\prelie\otimes\mathrm{id}) - \prelie\circ(\mathrm{id}\otimes\prelie)
\]
is left-symmetric, i.e., \(\mathrm{a}_{\prelie}=\mathrm{a}_{\prelie}\circ(\tau\otimes\mathrm{id})\),
and the following rectangle and all its inner squares commute:

\[
  \begin{tikzcd}
  T_1\otimes T_1\ar[r,"\partial^\otimes"]\ar[d,"\prelie"]&T_1\otimes T_0\oplus T_0\otimes T_1\ar[d,"\prelie"]\ar[r,"\partial^\otimes"]&T_0\otimes T_0\ar[d,"\prelie"]\\
  0\ar[r]&T_1\ar[r,"\partial"]&T_0
  \end{tikzcd}
\]
where \(\partial^\otimes(x\otimes y) = \partial(x)\otimes y+(-1)^{|x||y|}x\otimes\partial(y)\).

Unpacking this definition by restricting to each degree we see that, equivalently, this definition corresponds to a pre-Lie algebra \(\prelie\colon T_0\otimes T_0\to T_0\) together with a pre-Lie module structure \(\laction\colon T_0\otimes T_1\to T_1\), \(\raction\colon T_1\otimes T_0\to T_1\) and a map \(\partial\colon T_1\to T_0\) satsifying, for all \(s,t\in T_0\) and \(\mathscr{s},\mathscr{t}\in T_1\):
\begin{align}
  (s\prelie t-t\prelie s)\laction\mathscr{t} &= s\laction(t\laction\mathscr{t}) - t\laction(s\laction\mathscr{t})\tag{M1}\label{eq:m1}\\
  (s\laction \mathscr{t})\raction t - s\laction(\mathscr{t}\raction t) &= (\mathscr{t}\raction s)\raction t - \mathscr{t}\raction(s\prelie t)\tag{M2}\label{eq:m2},
\end{align}
and the equivariance and Peiffer relations
\begin{align}
  \partial(s\laction\mathscr{s}+\mathscr{t}\raction t) &= s\prelie\partial(\mathscr{s}) + \partial(\mathscr{t})\prelie t\tag{EQUI}\label{eq:equi}\\
  \partial(\mathscr{s})\laction\mathscr{t} &= \mathscr{s}\raction\partial(\mathscr{t})\tag{PEIF}\label{eq:peif}.
\end{align}

\begin{proposition}
Define the map \(\preliealt\colon T_1\otimes T_1\to T_1\) by
\[
  \mathscr{s}\preliealt\mathscr{t}\coloneqq\partial(\mathscr{s})\laction\mathscr{t}.
\]
Then \(T_1\) is a pre-Lie algebra and \(\partial\colon T_1\to T_0\) is a pre-Lie morphism.
\end{proposition}

We stress the fact that this second pre-Lie product is not immediately defined by the pre-Lie structure on the complex \((\partial\colon T_1\to T_0)\) but has to be defined a posteriori.

\begin{proof}
  To see that \(T_1\) is a pre-Lie algebra it suffices to note that by \cref{eq:m1} the associator
  \[
    \mathrm{a}_{\preliealt}(\mathscr{r},\mathscr{s},\mathscr{t}) = \mathrm{a}_{\prelie}(\partial(\mathscr{r}),\partial(\mathscr{s}),\mathscr{t})
  \]
  and thus it is left-symmetric.
  The second statement follows immediately from \cref{eq:equi}.
\end{proof}
\subsection{Pre-Lie crossed modules as internal categories}
Consider now the category \(\catpl\) of pre-Lie algebras with pre-Lie morphisms.
An internal category in \(\catpl\) is a pair \((C_0, C_1)\) of a pre-Lie algebra \((C_0,\prelie)\) of ``objects'' and a pre-Lie algebra \((C_1,\preliealt)\) of ``arrows'', maps \(s,t\colon C_1\to C_0\), the source and the target, an identity map \(i\colon C_0\to C_1\) and a composition \(\circ\colon C_1\times_{s,t} C_1\to C_1\) satisfying the standard axioms of a category, with all maps but the composition being pre-Lie morphisms.
The composition however satisfies the so-called \emph{interchange law}: whenever \(f,g,h,k\in C_1\) are composable
\[
  (f\circ g)\preliealt(h\circ k) = (f\preliealt h)\circ(g\preliealt k).
\]

\bibliographystyle{arxivalpha}
\bibliography{references}
\end{document}


